% Part: normal-modal-logic
% Chapter: axioms-systems
% Section: duals

\documentclass[../../../include/open-logic-section]{subfiles}

\begin{document}

\olfileid{nml}{prf}{dua}

\olsection{对偶\usetoken{P}{formula}}

\begin{defn}\ollabel{def:duals}
  公式 \Ax{T}、\Ax{B}、\Ax{4} 和 \Ax{5} 各有一个\emph{对偶}，用带下标的菱形符号表示如下：
  \begin{align}
    \tag{\Ax{T_\Diamond}} p & \lif \Diamond p\\
    \tag{\Ax{B_\Diamond}} \Diamond\Box p & \lif p\\
    \tag{\Ax{4_\Diamond}} \Diamond\Diamond p & \lif \Diamond p\\
    \tag{\Ax{5_\Diamond}} \Diamond\Box p & \lif \Box p
    \end{align}
\end{defn}

以上每个对偶公式都是从相应的原公式，通过将 $p$ 替换为 $\lnot p$、取逆否、将 $\lnot\Box\lnot$ 替换为 $\Diamond$，并将 $\lnot\Diamond\lnot$ 替换为~$\Box$ 得到的。\Ax{D}，即 $\Box!A \lif \Diamond!A$，在这个意义下是它自身的对偶。

\begin{prop}\ollabel{prop:dualsys}
  对于\olref[nml][prf][dua]{def:duals}中的每个公式~$!A$：$\Log{K}!A = \Log{K}!A_{\Diamond}$。
\end{prop}
\begin{proof}
  练习。
\end{proof}
\begin{prob}
  证明\olref[nml][prf][dua]{prop:dualsys}。
\end{prob}

\end{document}